% !TEX program = xelatex
% Arabic document showing both Arabic-Indic numerals (٠-٩)
% and Western numerals (0-9) side by side.
% Demonstrates how to switch between numeral systems.
\documentclass[12pt, a4paper]{article}

\usepackage{fontspec}
\usepackage{polyglossia}
\setmainlanguage[locale=mashriq]{arabic}
\setotherlanguage{english}

\newfontfamily\arabicfont[Script=Arabic]{Amiri}

\begin{document}

\title{الأرقامُ في النصِّ العربيّ}
\author{د. نبراس أبو الذهب}
\date{2025}
\maketitle

\section{الأرقامُ العربيةُ الهنديةُ والأرقامُ الغربية}

يَستخدمُ العربُ نظامينِ من الأرقامِ: الأرقامَ العربيةَ الهنديةَ
(المشرقية) والأرقامَ الغربيةَ (المغربية). يَتبيّنُ الفرقُ في الجدولِ أدناه.

\subsection{مقارنةُ الأرقام}

\begin{tabular}{|c|c|c|}
\hline
\LR{Western} & \LR{Arabic-Indic} & الرقم \\
\hline
0 & ٠ & صفر \\
1 & ١ & واحد \\
2 & ٢ & اثنان \\
3 & ٣ & ثلاثة \\
4 & ٤ & أربعة \\
5 & ٥ & خمسة \\
6 & ٦ & ستة \\
7 & ٧ & سبعة \\
8 & ٨ & ثمانية \\
9 & ٩ & تسعة \\
\hline
\end{tabular}

\subsection{التبديلُ بين النظامين}

لعرضِ الأرقامِ الغربيةِ داخلَ النصِّ العربيّ، نَستخدمُ الأمرَ \LR{\textbackslash LR\{\}}
مثل: \LR{2025} بدلاً من ٢٠٢٥.

كما يُمكنُ تَحديدُ نظامِ الأرقامِ عندَ إعدادِ اللغةِ عبرَ خيار
\LR{numerals=arabicdigits} أو \LR{numerals=maghrib} في
\LR{\textbackslash setmainlanguage}.

\end{document}
